\subsection{An executed coupling example on the deterministic-variance face}

At \(\xi=0,\kappa=3,\bar v=9/200,v_0=1/25,S_0=100,r=1/100,T=1,h=1/768\), retain the twelve
monthly fixings, the call spread with strikes 95 and 110, and
\(c=1254433/1250000\). Write
\(s_{\nu,j}=I_\nu(t_j)-I_\nu(t_{j-1})\), using the exact continuous and
Euler cumulative variances in (6.2). Couple the laws by twelve common
independent standard normals. With \(\sigma_\nu=\sqrt{s_{\nu,1}}\), remove the first
centered normal factor and put
\[
m_{\nu,i}=rt_i-I_\nu(t_i)/2,\qquad
\ell_{\nu,i,j}=\sqrt{s_{\nu,j}}\mathbf1_{\{j\le i\}},\quad 2\le j\le12,
\]
\[
a_{\nu,i}=100e^{m_{\nu,i}+\ell_{\nu,i}\cdot Z},\qquad
a_\nu=\frac1{12}\sum_i a_{\nu,i},\qquad
g_\nu=100e^{\bar m_\nu+\bar\ell_\nu\cdot Z},\qquad b_\nu=cg_\nu.
\]
Here \(Z=(Z_2,\ldots,Z_{12})\), and bars denote averaging over fixing
index \(i=1,\ldots,12\).
The drift includes the first-month term \(-s_{\nu,1}/2\); hence the removed
factor has the centered Gaussian law required by \(\mathsf H_s(c)\).
Here the stock driver \(\rho W+\sqrt{1-\rho^2}B\) is combined into a single
Brownian motion. On \(\xi=0\) the variance is deterministic and the price
law is independent of \(\rho\); hence \(\sigma_\nu^2=s_{\nu,1}\), using
the full first-month variance. The separate-law comparison below uses this
same conditioning.

All conditional second moments, including mixed P/Q moments, follow from
\[
E\!\left[100e^{m+u\cdot Z}\,100e^{n+v\cdot Z}\right]
=10000e^{m+n+\|u+v\|^2/2}.
\]
Define
\[
q_\nu=E(a_\nu-b_\nu)^2,\quad Z_2=q_P+q_Q,\quad
D_a=E(a_Q-a_P)^2,\quad D_b=E(b_Q-b_P)^2.
\]
For the deterministic \(M_K,N_K\) in the coupling proposition, the
pointwise bound \(z^2\le(a_P-b_P)^2+(a_Q-b_Q)^2\) and
Cauchy--Schwarz give the executable upper bound
\[
E\Xi_K\le B_K:=
M_K\sqrt{Z_2}(\sqrt{D_a}+\sqrt{D_b})
+\frac12N_KZ_2|\sigma_Q-\sigma_P|.
\]
The finite Gaussian-moment computation at 384 and 512 bits yields
\[
E\Xi_{95}\le .000023887739528734380066,\qquad
E\Xi_{110}\le .000020630320502088782784.
\]
Consequently the discounted nonlinear spread contribution belongs to the
symmetric interval with radius \(e^{-r}(B_{95}+B_{110})\), whose width is
at most \(.000088150195864703911227\). For the same laws, conditioning,
and scale, the separate-law estimate has width in
\[
[.015383637151510699801154,\;.015383637151510699801155].
\]
The ratio of widths is below \(.005731\). Strict improvement is checked
using the outward upper endpoint of the coupled width and the lower endpoint
of the separate width. The comparison uses
\(D_\nu=e^{2\sigma_\nu^2}q_\nu/\sigma_\nu\), so all weighted moments
refer to the same Gaussian factor.

This is a closed-moment enclosure of the coupled remainder expectations,
with no Monte Carlo or unbounded-domain quadrature. It establishes narrowing
of the nonlinear component at one deterministic-variance point. It does not
compute a complete price interval, a signed Asian leading coefficient, or a
coupled certificate for \(\xi>0\). The two precisions and two algebraic
Gaussian-moment expressions share the analytical proof and Arb; they are
arithmetic consistency checks.
